\section{Study area (Exercise 1.1)}\label{sec:study-area}
\noindent Responsible: Person A. Notebook: \code{notebooks/A1\_study\_area.ipynb}.

\subsection{Number of observations, persons and municipalities (1 point)}

\paragraph{Method.}
\todo{Name the chosen department and justify the choice. Describe how the detailed census was filtered to the department
(zone file, department code, \code{IRIS} / \code{CANTVILLE} columns). Explain the three counts: (1) observations are rows
of the filtered file, (2) persons are the \emph{sum of the weights} \code{IPONDI}, not the number of rows, (3) municipalities
are counted from a source that identifies all of them. State how the municipalities without IRIS code in the detailed census
(small municipalities, localised only at ``canton-ou-ville'' level) are handled.}

\paragraph{Result.}
\todo{Report the three numbers, ideally in \cref{tab:study-area-counts}, and comment on the difference between observations
and weighted persons.}

\begin{table}[H]
  \centering
  \caption{\todo{Caption: census counts for the study area.}}
  \label{tab:study-area-counts}
  \begin{tabular}{@{}lr@{}}
    \toprule
    Quantity & Value \\
    \midrule
    Person observations (rows, unweighted) & \todo{value} \\
    Persons (sum of weights \code{IPONDI}) & \todo{value} \\
    Municipalities & \todo{value} \\
    \bottomrule
  \end{tabular}
\end{table}

\subsection{Map of the department inside of France (1 point)}

\paragraph{Method.}
\todo{Describe how the map is built from the IGN CONTOURS-IRIS 2021 data (dissolve to departments, highlight the study area,
projection EPSG:2154).}

\begin{figure}[H]
  \centering
  \figplaceholder{study_area_location}
  \caption{\todo{Caption: location of the study area inside of France.}}
  \label{fig:study-area-location}
\end{figure}

\paragraph{Result.}
\todo{Describe what \cref{fig:study-area-location} shows (neighbouring departments, main cities).}

\paragraph{Problems encountered.}
\todo{Document every problem for partial points: e.g.\ census files per zone, missing IRIS codes, municipality codes that
differ between sources, how they were resolved.}
